% Conversão fiel de chapter-02-text-h-0009; origem RELATORIO_ATUAL.odt.
\hypertarget{__RefHeading___Toc4735_3118176633}{}
\subsection{Banco de dados utilizado\originalspace{}}
\label{ages-i:text-h-0009}

\indent \originaltab{}O banco de dados utilizado é do tipo relacional com base em um modelo do \originalinlinetab{}Postgresql\space\cite{postgresql-docs}. Também foi utilizado o Prisma\space\cite{prisma-docs} como ferramenta de ORM para \originalinlinetab{}facilitar o acesso e manipulação dos dados na aplicação. Essa combinação \originalinlinetab{}oferece um ambiente robusto e eficiente para o armazenamento e \originalinlinetab{}gerenciamento de dados.\par

\noindent \originaltab{}Mais detalhes podem ser encontrados aqui: \href{https://tools.ages.pucrs.br/easy-choose/wiki/-/wikis/Banco\%20de\%20Dados}{Banco de Dados}\par

\memorialfigure{imagens/user-supplied/ages-i-database.png}{\textbf{Figura 2 –} Diagrama Banco de Dados - AGES I}{fig:ages-i-database}{Fonte: \url{https://tools.ages.pucrs.br/easy-choose/wiki/-/wikis/Banco-de-Dados}}{0.92}
